\chapter{运动学基础}
\label{sec:02-fundamentals-of-kinematic}

\section{机器人的自由度}

自由度 (Degrees of freedom, DOF)：确定
机构位形所需独立参数的数目。

\section{描述：位置、姿态与坐标系}

$$
\text{位姿（Location）} = 
\begin{array}[t]{c}
\text{位置（Position）} \\
\Downarrow \\
3
\end{array}
+
\begin{array}[t]{c}
\text{姿态（Orientation）} \\
\Downarrow \\
3
\end{array}
$$

位置描述
$$
{ }^A P=\left(\begin{array}{l}
p_x \\
p_y \\
p_z
\end{array}\right)
$$

姿态描述：旋转矩阵
$$
{ }_B^A R=\left[{ }^A \hat{X}_B{ }^A \hat{Y}_B{ }^A \hat{Z}_B\right]=\left[\begin{array}{lll}
r_{11} & r_{12} & r_{13} \\
r_{21} & r_{22} & r_{23} \\
r_{31} & r_{31} & r_{33}
\end{array}\right]
$$

对应于轴$x$，$y$或$z$作转角为$\theta$的旋转变换，其旋转矩阵分别为：
$$
R(x, \theta)=\left[\begin{array}{ccc}
1 & 0 & 0 \\
0 & \cos \theta & -\sin \theta \\
0 & \sin \theta & \cos \theta
\end{array}\right]
$$
$$
R(y, \theta)=\left[\begin{array}{ccc}
\cos \theta & 0 & \sin \theta \\
0 & 1 & 0 \\
-\sin \theta & 0 & \cos \theta
\end{array}\right]
$$
$$
R(z, \theta)=\left[\begin{array}{ccc}
\cos \theta & -\sin \theta & 0 \\
\sin \theta & \cos \theta & 0 \\
0 & 0 & 1
\end{array}\right]
$$

位姿（Location）＝位置（Position）＋姿态（Orientation）；位姿描述 = 位置矢量 + 旋转矩阵。

表示位置时，旋转矩阵等于单位矩阵；表示姿态时，位置矢量等于零。

$$
\{B\} = \{ \underbrace{{}^{A}_{B}R}_{\text{旋转矩阵}}, \quad \underbrace{{}^{A}P_{BORG}}_{\text{原点位置矢量}} \}
$$

手爪的方位——旋转矩阵$R=\left[\begin{array}{lll}n & o & a\end{array}\right]$，描述手爪的姿态。

$$
\{T\}=\left\{\begin{array}{llll}
n & o & a & p
\end{array}\right\}=
\left(\begin{array}{cccc}
n_1 & o_1 & a_1 & \tikzmarknode{x}{x} \\
n_2 & o_2 & a_2 & \tikzmarknode{y}{y} \\
n_3 & o_3 & a_3 & \tikzmarknode{z}{z} \\
0 & 0 & 0 & 1
\end{array}\right)
$$

\begin{tikzpicture}[overlay, remember picture]
    \node[draw=whusakura, thick, ellipse, inner xsep=2pt, inner ysep=0pt, fit=(x) (y) (z)] (circle) {};
    \draw[<-, whusakura, thick] (circle.east) -- ++(1,0) node[right] {坐标原点位置};
\end{tikzpicture}

\section{映射：从坐标系到坐标系的变换}

\subsection{平移映射}

\(\{B\}\)与\(\{A\}\)方位（姿态）相同，但原点不重合，用位置矢量${ }^A P_{B_0}$表示\(\{B\}\)的原点相对于\(\{A\}\)的位置。

坐标平移方程：
$$
{ }^A p = { }^B p + { }^A p_{B_0}
$$

\subsection{旋转映射}

坐标旋转方程
$$
{ }^A P = { }^A_B R { }^B P
$$

旋转的叠加：绕固定轴旋转（前乘） / 绕当前轴旋转（后乘）

\subsection{一般坐标系的映射}

$$
{ }^{A} p = \tikzmarknode{R}{{}^{A}_{B} R} \, { }^{B} p + \tikzmarknode{P}{{ }^{A} p_{B_0}}
$$

% 在后台绘制箭头和文字
\begin{tikzpicture}[overlay, remember picture, >=stealth]
    % 向下引出标注旋转矩阵
    \draw[<-, thick, whusakura] (R.south) -- ++(0,-0.6)
        node[below, align=center, font=\small] {$\{B\}$ 相对于 $\{A\}$ 的 \\ (方位)旋转矩阵};
    % 向上引出标注平移向量
    \draw[<-, thick, whusakura] (P.east) -- ++(0.6,0)
        node[right, align=center, font=\small] {$\{B\}$ 原点相对于 \\ $\{A\}$ 的位置};
\end{tikzpicture}

\vspace{0.4in}

为了简化计算，将上述包含加法的仿射变换统一为单一的线性矩阵乘法，引入齐次坐标变换。定义$4 \times 4$ 的齐次变换矩阵 ${}^{A}_{B}T$：
$$
    {}^{A}_{B}T = 
    \begin{bmatrix}
        {}^{A}_{B}R & {}^{A}p_{B_0} \\
        0 & 1
    \end{bmatrix}
$$

\vspace{.6in}

\begin{equation*}
    % 使用 \tikzmarknode 为齐次坐标向量、旋转矩阵、平移矢量分别打上标签
    \tikzmarknode{vecA}{\begin{bmatrix} {}^{A} p \\ 1 \end{bmatrix}} 
    = 
    \begin{bmatrix}
        \tikzmarknode{matR}{{}^{A}_{B}R} & \tikzmarknode{vecP}{{}^{A}p_{B_o}} \\
        0 & 1
    \end{bmatrix}
    \tikzmarknode{vecB}{\begin{bmatrix} {}^{B}p \\ 1 \end{bmatrix}}
\end{equation*}

% 在后台图层利用打好的标签绘制红圈、虚线箭头和文字注释
\begin{tikzpicture}[remember picture, overlay, >=stealth]
    \node[draw=red, dashed, ellipse, inner xsep=4pt, inner ysep=8pt, fit=(vecA)] (ellA) {};
    \node[draw=red, dashed, ellipse, inner xsep=4pt, inner ysep=8pt, fit=(vecB)] (ellB) {};

    \node[above=.1in of ellA] {齐次坐标};
    \node[above=.1in of ellB] {齐次坐标};

    \node[below=.6in of matR, xshift=-.2in] (labelR) {旋转矩阵};
    \node[below=.6in of vecP, xshift=.2in] (labelP) {平移矢量};
    
    % 从矩阵元素引出红色虚线箭头指向底部文字
    \draw[->, red, dashed, thick] (matR.south) -- (labelR.north);
    \draw[->, red, dashed, thick] (vecP.south) -- (labelP.north);
\end{tikzpicture}

\vspace{.6in}

% 最终简化的齐次变换紧凑形式
$$
    {}^{A} P = {}^{A}_{B} T \, {}^{B} P
$$
